
\documentclass[11pt]{article}
\usepackage[margin=1in]{geometry}

% Core packages
\usepackage{amsmath,amssymb}
\usepackage{tikz-cd}
\usepackage{multicol}
\usepackage{hyperref}

% Paragraphs
\setlength{\parindent}{0pt}
\setlength{\parskip}{1\baselineskip}

\title{Boundary Differential Equation Visualization}
\author{Nam Do\\University of Oulu}

\date{\today}

\begin{document}
\maketitle
\section{Autonomous ODE and Vector Field}
The foundation of continuos time dynamics is the ordinary differential equation (ODE), we work exclusively with the autonomous case, where the rule governing the system's evolution does not depend explicitly on time.\\
\textbf{Definition 1.1 (Autonomous ODE)}
Let \(f: \mathbb{R}^d \to \mathbb{R}^d \) be a \(C^1\) vector field. An \textbf{autonomous ordinary differential equation} on \(\mathbb{R}^d\) is an equation of the form
\[
    \dot{x}(t) = f(x(t)), \qquad x(t_0) = x_0 \in \mathbb{R}^d. \tag{1}
\]
where \(\dot{x}(t) = dx/dt\) denotes the time derivative of the state. The solution \(t \mapsto x(t)\) is called \textbf{trajectory} or \textbf{integral curve} of $f$.
\section{Random ODE with Bounded Noise}
We are interested in continuous dynamical systems' behavior when some noise is added in each state of the system. We do not care about the probability distribution of the noise, but we have a hard constraint of the noise amplitude (i.e, the noise is bounded). \\
\textbf{Defintion 2.1 (Random differential equation with bounded noise).} 
Let \( f : \mathbb{R}^d \to \mathbb{R}^d \) be a \( C^1 \) vector field and \(\epsilon > 0\). A \textbf{random differential equation with bounded noise} is
\[
    \dot{x}(t) = f(x(t)) + \xi_t, \quad \xi_t \in B_\epsilon(0) \subset \mathbb{R}^d, \tag{2}\label{eq:rde}
\]

where \(\xi_t \in B_\epsilon(0)\) is the velocity correction at time $t$ with the maginitude at most $\epsilon$, centered at the zero vector. More formally, the noise signal $\xi: \mathbb{R} \to B_\epsilon(0)$ ranges over the class
\[
\mathcal{U} = \bigl\{\, \xi : \mathbb{R} \to B_\varepsilon(0) \;\big|\; \xi \text{ measurable} \,\bigr\}, \tag{3}
\]
of all measurable bounded-noise functions, and \(B_\epsilon(0) = \{v \in \mathbb{R}^d : \|v\| \leq \varepsilon \}\) is the close ball of radius $\epsilon$ centered at the origin.

\section{Set-valued Reformulation}
The RDE~\eqref{eq:rde} ioerates on individual trajectories: for a fixed noise singal $\xi \in \mathcal{U} $ and a fixed initial condition $x_0$, it produces a unique solution curve \(t \mapsto \phi_t(x_0;\xi)\). But since we make no assumption about which $\xi$ is realised, we must condiere the entire family of trajectories simultaneously. This shifts from tracking a single point to tracking the set of all possible point, so now we derive the set-valued random ODE definition\\
\textbf{Definition 3.1 (Set-valued map with bounded noise)}. Let $f: \mathbb{R}^d \to \mathbb{R}^d$ be a $C^1$ vector field and $\epsilon > 0$. The \textbf{set-valued map} associated with $f$ and noise radius $\epsilon$ is
\[
    F_\epsilon: \mathcal{K}(\mathbb{R}^d) \to \mathcal{K}(\mathbb{R}^d), \quad F_\epsilon(M) := \bigcup_{x \in M} B_\varepsilon\!\bigl(f(x)\bigr), \tag{4}
\]
where $K(\mathbb{R}^d)$ denotes the collection of all compact subsets of $\mathbb{R}^d$
\section{Boundary tracking of set-valued systems}
Consider a boundary point $x \in \partial M$ with outward unit normal \(\hat{n}\). Given that the system is at $x \in \partial M$, we want to find out noise direction $\xi \in B_\epsilon(0)$ is worst-case, i.e., which $\xi$ pushes the next state as far as possibe toward the exeterior of M. 

The worst-case next - the one lying on the boundary of $F_{\epsilon}(\{x\}) = B_\epsilon(f(x))$ that is also on $\partial M$ is
\[
    y = f(x) + \epsilon\hat{n}, \tag{5}\label{eq:single_boundary}
\]
where $\hat{n}'$ is the outward unit normal to $\partial M$ at $y$. 

\subsubsection{The boundary map}
The observation in ~\eqref{eq:single_boundary} defines a single-valued map on the sace $\mathbb{R}^d \times S^{d-1}$ of (position, outward normal) pairs. A point on $\partial M$ carries not just its position $x$ but also its outward unit normal $\hat{n}$; together $(x,\hat{n})$ is a point in the \textbf{outward unit normal bundle} $UN^{+}\partial M$. Following the approach developed by Tey~\cite{tey2023}, we define the boundary map as follows:

\textbf{Definition 6 (Boundary map)} \cite{tey2023}: The \textbf{boundary map} $E: \mathbb{R}^d \times S^{d-1} \to {R}^d \times S^{d-1}$ is defined by

\[
E(x,\hat{n})=\left(f(x)+\varepsilon\,\frac{n'}{\|n'\|},\;\frac{n'}{\|n'\|}\right),
\qquad
n'=\bigl(Df(x)^{-\top}\bigr)\hat{n}, \tag{6}
\]

where $Df(x)^{-\top} = \bigl(Df(x)^{-1}\bigr)^\top$ is the inverse-transpose of the Jacobian of \(f\) at \(x\).
The two components of E encode exactly the two steps of worst-case boundary evolution: \\
\textbf{Position update} \(x \mapsto f(x) + \varepsilon\,\hat{n}'\): apply the deterministic map $f$, then add the worst-case noise $\epsilon\hat{n}'$ in the outward normal direction. This places the new boundary point at the outermost edge of the $\epsilon$-ball around $f(x)$\\ 
\textbf{Normal update} \(n \mapsto n'/\|n'\|\) with \(n' = Df(x)^{-\top}n\): propagate the outward normal through the non-linear map $f$ using the inverse-transpose of the Jacobian.
\section{Discretization of Random ODE with bounded noise}
We use RK4 method to discretize continuous dynamical systems for visualization. To approximate the next system with respect to time, first, we used to Euler method to discretize the system state, and then refining with the RK4 method for better approximation. In this case, we use the explicit Euler method with step size $h > 0$ is given by:
\[
    x_{i+1} = x_i + hf(x_i) + h\xi_{t_{i}}, \quad h\xi_{t_{i}} \in B_{h\epsilon}(0)
\]


This is exactly a discrete set-valued map with
\[
\hat{f}(x)=x+h f(x), \qquad \varepsilon' = h\varepsilon. \tag{7}
\]
The set-valued mapping corresponding to the Euler step is therefore
\[
\hat{F}_{h\varepsilon}(x)=B_{h\varepsilon}\!\bigl(\hat{f}(x)\bigr) = B_{h\varepsilon}\!\bigl(x+h f(x)\bigr), \tag{8}
\]
which has exactly the form \(F_{\varepsilon'}(x)=B_{\varepsilon'}(f(x))\) studied throughout your project.  
The MFI set of the continuous system (10) is approximated by the MIS of the discrete system with map \(\hat{f}\) and noise \(\varepsilon' = h\varepsilon\), with the approximation improving as \(h \to 0\).
\subsection{Runge--Kutta Methods}
The standard ODE integrator is \textbf{RK4}:
\[
\begin{aligned}
k_1 &= g(z_i),\\
k_2 &= g\!\left(z_i+\frac{h}{2}k_1\right),\\
k_3 &= g\!\left(z_i+\frac{h}{2}k_2\right),\\
k_4 &= g(z_i+h k_3),\\
z_{i+1} &= z_i+\frac{h}{6}(k_1+2k_2+2k_3+k_4).
\end{aligned} \tag{9}
\]

\subsection{The Euler Boundary Map as a Discrete Scheme}
The Euler boundary map is a numerical scheme for the BDE. Its advantages:
\begin{enumerate}
\item It maps exactly onto your existing discrete boundary map code:
\[
\hat x = f + h\cdot \hat f
\]
as the underlying map, \(\varepsilon' = h\varepsilon\).
\item No separate ODE solver is needed---reuse Algorithm 1 (boundary map evolution) from~\cite{tey2023}.
\item The approximation order is \(O(h)\) (Euler), so small \(h\) is required.
\end{enumerate}
\paragraph{Algorithm 1: RK4 Boundary Map Step}
\textbf{Require:} State \((x,\hat n)\in \mathbb{R}^d\times S^{d-1}\), vector field \(f\), Jacobian \(Df\), noise \(\varepsilon\), step \(h\).

Define
\[
\Phi(x,\hat n)=
\begin{bmatrix}
f(x)+\varepsilon \hat n\\[2mm]
\dfrac{Df(x)^{-\top}\hat n}{\left\|Df(x)^{-\top}\hat n\right\|}
\end{bmatrix}.
\]

Then RK4 on \((x,\hat n)\):
\[
\begin{aligned}
(K_1^x,K_1^n)&=\Phi(x_i,\hat n_i),\\
(K_2^x,K_2^n)&=\Phi\!\left(x_i+\frac{h}{2}K_1^x,\ \hat n_i+\frac{h}{2}K_1^n\right),\\
(K_3^x,K_3^n)&=\Phi\!\left(x_i+\frac{h}{2}K_2^x,\ \hat n_i+\frac{h}{2}K_2^n\right),\\
(K_4^x,K_4^n)&=\Phi\!\left(x_i+hK_3^x,\ \hat n_i+hK_3^n\right),\\[1mm]
x'&=x_i+\frac{h}{6}\left(K_1^x+2K_2^x+2K_3^x+K_4^x\right),\\
\hat n'&=\hat n_i+\frac{h}{6}\left(K_1^n+2K_2^n+2K_3^n+K_4^n\right),\\
\hat n'&\leftarrow \frac{\hat n'}{\|\hat n'\|}.
\end{aligned}
\]
Return \((x',\hat n')\).

\begin{thebibliography}{9}
\bibitem{tey2023}
Wei Hao Tey.
\textit{On Minimal Invariant Sets of Certain Set-Valued Dynamical Systems -- A Boundary Map Approach}.
PhD thesis, Imperial College London, December 2022. Available at: \url{https://www.researchgate.net/publication/366016533_On_Minimal_Invariant_Sets_of_Certain_Set-Valued_Dynamical_Systems_-_A_Boundary_Map_Approach}
\end{thebibliography}

\end{document}

